% Appliquer la relation entre $\Delta E_c$ et $\sum W\left(\vec{F}_{\text {ext }}\right)$

Une caisse, de masse $m=\qty{25.0}{\kg}$ est amenée en haut d'un plan incliné de
l'angle $\alpha=\ang{15}$ par rapport à l'horizontale.

Le plan incliné est supposé parfaitement lisse.

On lâche la caisse sans vitesse initiale et elle descend le long de la ligne de
plus grande pente du plan incliné. Calculer la valeur de la vitesse de la caisse
lorsqu'elle a parcouru une distance $\ell=\qty{5.00}{\m}$ suivant cette ligne.

\begin{donnees}
    Intensité de la pesanteur~: $g=\qty{9.8}{\N\per\kg}$.
\end{donnees}
